\section{Experimental Design}

\subsection{Paired common-support estimand}

Let $y_{i,t}$ denote the realized demand for series $i$ at target time $t$. Let $s_i$ be the grounded history start and $f_i$ the first-positive time. For forecasting method $m$, the two one-step forecasts are
\begin{align}
\hat y^{G}_{i,t,m} &= \mathcal{F}_m(y_{i,s_i:t-1}),\\
\hat y^{T}_{i,t,m} &= \mathcal{F}_m(y_{i,f_i:t-1}),
\end{align}
where $G$ denotes grounded/lifecycle-aligned history and $T$ denotes first-positive trimming. The target $y_{i,t}$ is identical in both conditions.

For absolute error, we first average the paired error difference within each series over its common-support target set $\mathcal{T}_i$, and then average equally across the $N$ affected series:
\begin{equation}
\Delta_m =
\frac{1}{N}
\sum_{i=1}^{N}
\left[
\frac{1}{|\mathcal{T}_i|}
\sum_{t\in\mathcal{T}_i}
\left(
|y_{i,t}-\hat y^T_{i,t,m}|
-
|y_{i,t}-\hat y^G_{i,t,m}|
\right)
\right].
\end{equation}
This is a series-weighted estimand: each product--store series receives equal weight regardless of how many common-support origins it contributes. Positive $\Delta_m$ means trimming worsens MAE; negative values mean trimming improves MAE.

The primary minimum-history requirement is three observations in both conditions. Visuelle additionally checks minimum histories of 2 and 4. FreshRetailNet-50K checks 2, 3, 4, and 7 observations. FreshRetailNet-50K uses prespecified target days
\[
\{4,7,14,21,28,42,56,70,84,89\},
\]
which allows the history-start effect to be examined from very short to long histories.

\subsection{Forecasting methods}

We use nine canonical implementations from StatsForecast 2.1.1 \citep{garza2022statsforecast}:
Naive, HistoricAverage, optimized simple exponential smoothing (SESOpt), CrostonClassic, CrostonOptimized, CrostonSBA, TSB, ADIDA, and IMAPA. TSB uses $\alpha_d=\alpha_p=0.2$ in the primary analysis; Visuelle additionally checks 0.1 and 0.3. Forecasts are truncated below at zero where applicable.

These methods span last-observation, mean-based, exponential-smoothing, Croston-style, probability-update, and temporal-aggregation families. The purpose is not to identify a state-of-the-art winner, but to test whether sensitivity generalizes across forecasting mechanisms.

\subsection{Uncertainty and rank stability}

For each method we first average forecast errors within each store--product series across eligible origins. We then use a product-cluster bootstrap with 5,000 replicates. Products are sampled with replacement as clusters, but all store-series belonging to sampled products are pooled, preserving the same \emph{series-weighted} estimand as the reported point estimate.

For each replicate we compute $\Delta_m$, model ranks under both history conventions, Kendall's $\tau$, and the winning method. This corrected procedure is important: a preliminary bootstrap that averaged products equally targeted a different product-weighted estimand. The final results reported here use cluster resampling with series weighting throughout.

\subsection{Eligibility and global propagation}

The common-support experiment deliberately conditions on targets forecastable under both histories. Trimming can also delay the earliest forecast origin. For a minimum history $h$, we define
\begin{equation}
\text{Coverage loss}
=
1-\frac{N_T(h)}{N_G(h)},
\end{equation}
where $N_G$ and $N_T$ are the numbers of eligible forecast origins under the grounded and trimmed histories.

Finally, we propagate the history convention to the full benchmark. Unaffected series have identical forecasts under both conditions; affected series use their alternative histories. Visuelle global evaluation uses target weeks 8--11. FreshRetailNet-50K uses target days 28, 56, and 89. We report global MAE, WAPE, and deterministic rank changes. FreshRetailNet-50K also includes a sensitivity analysis that trims all 3,256 delayed-first-positive series regardless of stock confirmation; this ``blanket'' scenario is operational rather than causally grounded.

\subsection{Scaled metrics}

MASE scales absolute forecast error by an in-sample naive-error denominator \citep{hyndman2006mase}. We report two variants:
\begin{itemize}[leftmargin=1.5em]
    \item \textbf{Fixed denominator}: both forecast conditions use the grounded-history scale. This isolates forecast changes.
    \item \textbf{Convention-specific denominator}: each condition computes its scale from its own history. This captures forecast changes plus metric-normalization changes.
\end{itemize}
We apply the same distinction to RMSSE. This separation prevents a history-start intervention from silently changing both the prediction and the measurement scale.
